\section{Proof of the main theorem}\label{sec:conclusion}

We assemble the component arguments, retaining the distinction between
top-degree homology at odd primes and propagation at two.

\begin{proof}[Proof of Theorem~\ref{thm:main}]
For odd $p$, Proposition~\ref{prop:unitary-odd} proves the dimension
property for every unitary extension required by
Theorem~\ref{thm:reduction-odd}. That reduction proves the theorem.

Suppose now that $p=2$ and that a counterexample exists. Choose one of
minimal order. By Theorem~\ref{thm:reduction-two} it has a simple
component in the displayed classical list.
Proposition~\ref{prop:linear-two} excludes the linear and unitary
groups of natural dimension at least five, and
Proposition~\ref{prop:symplectic} excludes the symplectic groups on
the list. The standard isomorphisms
\[
\PSL_4(q)\cong\mathrm P\Omega_6^+(q),
\qquad
\PSU_4(q)\cong\mathrm P\Omega_6^-(q)
\]
place the remaining linear and unitary groups among the orthogonal
cases. Proposition~\ref{prop:orthogonal} excludes all these
orthogonal groups except the square-discriminant six-dimensional
case over $\F_5$. That case is $\PSL_4(5)$ and is excluded by
Proposition~\ref{prop:psl45}. The list is exhausted, a contradiction.
\end{proof}
